\par\noindent\textbf{(1)} 对 g(x) 求导 g'=RT ln[x/(1−x)]+Ω(1−2x)。组分化学势满足 μ\_A−μ\_B=g'，而局域混合稳定要求 g''=RT/[x(1−x)]−2Ω>0。若二阶曲率为负，任意小的成分涨落可降低自由能。
\par\noindent\textbf{(2)} 自旋线由 g''=0 给 T\_s(x)=2Ωx(1−x)/R，最大值在 x\_c=1/2，对应 T\_c=Ω/(2R)。因模型关于 x↔1−x 对称，临界点还满足 g'''=0；高于T\_c所有x的 g''>0，无相分离自旋不稳定区。
\par\noindent\textbf{(3)} 二相共存组成 x\_α,x\_β 应满足 g'(x\_α)=g'(x\_β)=[g(x\_β)−g(x\_α)]/(x\_β−x\_α)，即公共切线两端共用化学势。对对称模型共存端满足 x\_β=1−x\_α 且 g'(x\_α)=0；自旋线仅给局部不稳定边界，不能当成共存双相线。
\par\noindent\textbf{(4)} 在自旋线内 g''<0，浓度扰动可无势垒增长，称自旋分解；共存线与自旋线之间虽二相态更稳定，局部自由能仍凸，形成新相通常需成核势垒。临界附近有限梯度能和界面张力决定实际微观动力学，零梯度混合模型只给均匀态热力学。
